% Modelo UNISINOS para teses e dissertacoes

% PARA COMEÇAR A ESCREVER NÃO SE DETENHA AOS CÓDIGOS INICIAIS DO DOCUMENTO, VOCÊ IRÁ ENTENDÊ-LOS COM O TEMPO. 
% vÁ DIRETO PARA O CAPÍTULO 1 - INTRODUÇÃO.

% Arquivo de configurações e pacotes.
\input{0-setup}

% Arquivo de dados do documento: título, autor...
\input{1-dados}

% ----------------------------------------------
% Início do documento
% ----------------------------------------------
\begin{document}

% ----------------------------------------------
% Adiciona lista de correções no início do documento.
% Comentar a linha abaixo quando o trabalho for concluído
% ----------------------------------------------
\listoftodos
% ----------------------------------------------

% Arquivo de elementos pré-textuais: capa, folha de rosto, ficha catalografica, errata, folha de aprovação dedicatória, agradecimentos, epígrafe, resumos, lista de abreviaturas e siglas, lista de símbolos... 
\input{2-pretextual}

% ----------------------------------------------
% ELEMENTOS TEXTUAIS
% ----------------------------------------------
\textual

% ----------------------------------------------
% Exemplo de capítulo sem numeração, mas presente no Sumário
% ----------------------------------------------
%\chapter*[Introdução]{Introdução}
%\addcontentsline{toc}{chapter}{Introdução}
% ----------------------------------------------

% ||||||||||||||||||||||||||||||||||||||||||||||
% CAP 1 - INTRODUÇÃO 
% ||||||||||||||||||||||||||||||||||||||||||||||
\chapter[Introdução]{Introdução}

%\todo[inline, color=yellow]{exemplo de comentário para auxiliar na lista de tarefas e correções}

\todo[inline]{Alguns manuais de pacotes latex foram adicionados na pasta manuals}

\todo[inline, color=red]{Procure se informar a respeito do Mendeley. É um software para gerenciar referencias bibliográficas que irá facilitar muito a inclusão das referências no documento. Procure também como importar as referencias inseridas no mendeley para o arquivo references.bib aqui no overleaf}

Este documento e seu código-fonte são exemplos de referência de uso da classe \textsf{abntex2} e do pacote \textsf{abntex2cite}. O documento exemplifica a elaboração de trabalho acadêmico (tese, dissertação e outros do gênero) produzido conforme a ABNT NBR 14724:2011 \emph{Informação e documentação - Trabalhos acadêmicos - Apresentação}.

A expressão ``Modelo Canônico'' é utilizada para indicar que \abnTeX\ não é modelo específico de nenhuma universidade ou instituição, mas que implementa tão somente os requisitos das normas da ABNT. Uma lista completa das normas observadas pelo \abnTeX\ é apresentada em \citeonline{abntex2classe}. 

Este documento deve ser utilizado como complemento dos manuais do \abnTeX\ \cite{abntex2classe,abntex2cite,abntex2cite-alf} e da classe \textsf{memoir} \cite{memoir}. 

Esperamos, sinceramente, que o \abnTeX\ aprimore a qualidade do trabalho que você produzirá, de modo que o principal esforço seja concentrado no principal: na contribuição científica.

\section{Justificativa}
	\lipsum[1]

\section{Delimitação do tema ou problema}
	\lipsum[2]

\section{Delimitação do trabalho}

\section{Objetivos}

\subsection{Objetivos Gerais}
	\lipsum[3]
    
\subsection{Objetivos Específicos}
	\lipsum[4]
    

% ||||||||||||||||||||||||||||||||||||||||||||||
% CAP 2 - REVISÃO BIBLIOGRAFICA
% ||||||||||||||||||||||||||||||||||||||||||||||
\chapter{Revisão Bibliográfica}

\begin{figure}[htp]
	\centering
	\caption{\label{fig:met-disc-fig01} No centro da figura é representado um logo $U_1$.} 
	\includegraphics[width = 0.8\linewidth]{images/unisinos.png}
	\legend{Fonte: Desenvolvido pelo autor.}
\end{figure}

Matematicamente, o Fator de Potência (FP) pode ser expresso como:
\begin{equation}
	\label{eq:k-55}
    {
    \displaystyle 
    FP = \frac{\cos(\varphi)}{\sqrt{1 - THD^2}}
    }
\end{equation}

\section{Desenho de circuitos}

\begin{circuitikz} \draw
	(0,0)
    %to [battery, v=V1]
    to [R, l=R1, i=$i_{1}$]
    (0,4)
    
    (0,4)
    to [battery, v=V1]
    %to [R, l=R1, i=$i_{1}$]
	(4,4)
    node[anchor=south]{A}
    
    (4,0)
    to [R, l=R2, v>=$V_{R2}$, i=$i_{2}$, *-*]
    (4,4)
    
    (4,0)
    to [short]
    (0,0)
    
    (8,4)
    to [battery, v_>=V2]
    %to [R, l=R3, i>^=$i_{3}$]
	(4,4)
    
    (8,0)
    to [short]
    (4,0)
    node[anchor=north]{G}
    
    (8,0)
    %to [battery, v=V2]
    to [R, l=R3, i=$i_{3}$]
    (8,4)
    ;
\end{circuitikz}


\section{Outra linha de ref}

% ----------------------------------------------
%\section{Citações}
% ----------------------------------------------

%\index{citações!diretas}Utilize o ambiente \texttt{citacao} para incluir citações diretas com mais de três linhas:

%\begin{citacao}
%As citações diretas, no texto, com mais de três linhas, devem ser destacadas com recuo de 4 cm da margem esquerda, com letra menor que a do texto utilizado e sem as aspas. No caso de documentos datilografados, deve-se observar apenas o recuo \cite[5.3]{NBR10520:2002}.
%\end{citacao}

%\index{citações!simples}Citações simples, com até três linhas, devem ser incluídas com aspas. Observe que em \LaTeX~ as aspas iniciais são diferentes das finais: ``Amor é fogo que arde sem se ver''. 


% ----------------------------------------------
%\section{Remissões internas}
% ----------------------------------------------

%Ao nomear a \autoref{tab-nivinv}, apresentamos um exemplo de remissão interna, que também pode ser feita quando indicamos o \autoref{introdução}\footnote{O número do capítulo indicado é \ref{cap_exemplos}, que se inicia à página \pageref{cap_exemplos}.} (\nameref{cap_exemplos}, \autopageref{cap_exemplos}), por exemplo.

% ----------------------------------------------
%\section{Tabelas}
% ----------------------------------------------

%Apresenta-se um exemplo de tabela a ser confeccionada. Atente-se para as normas de tabela exigidas pela Universidade.

%\index{tabelas}A \autoref{tab-nivinv} é um exemplo de tabela construída em \LaTeX.

%\begin{table}[htb]
%\footnotesize
%\caption[Níveis de investigação]{Níveis de investigação.}
%\label{tab-nivinv}
%\begin{tabular}{p{2.6cm}|p{6.0cm}|p{2.25cm}|p{3.40cm}}
  %\hline
%   \textbf{Nível de Investigação} & \textbf{Insumos}  & \textbf{Sistemas de Investigação}  & \textbf{Produtos}  \\
%    \hline
%    Meta-nível & Filosofia\index{filosofia} da Ciência  & Epistemologia &
%    Paradigma  \\
%    \hline
%    Nível do objeto & Paradigmas do metanível e evidências do nível inferior &
%    Ciência  & Teorias e modelos \\
%    \hline
%    Nível inferior & Modelos e métodos do nível do objeto e problemas do nível inferior & Prática & Solução de problemas  \\
   % \hline
%\end{tabular}
%\legend{Fonte: \citeonline{van86}}
%\end{table}

% ----------------------------------------------
%\section{Expressões matemáticas}
% ----------------------------------------------

%\index{expressões matemáticas}Use o ambiente \texttt{equation} para escrever expressões matemáticas numeradas:

%\begin{equation}
%  \forall x \in X, \quad \exists \: y \leq \epsilon
%\end{equation}

%Escreva expressões matemáticas entre \$ e \$, como em $ \lim_{x \to \infty} \exp(-x) = 0 $, para que fiquem na mesma linha.

%Também é possível usar colchetes para indicar o início de uma expressão matemática que não é numerada.

%\[
%\left|\sum_{i=1}^n a_ib_i\right|
%\le
%\left(\sum_{i=1}^n a_i^2\right)^{1/2}
%\left(\sum_{i=1}^n b_i^2\right)^{1/2}
%\]

%Consulte mais informações sobre expressões matemáticas em \url{http://code.google.com/p/abntex2/w/edit/Referencias}.

%\section{Figuras}

%\index{figuras}Figuras podem ser criadas diretamente em \LaTeX, como o exemplo da \autoref{fig_circulo}.

%\begin{figure}[htb]
%	\begin{center}
%	    \setlength{\unitlength}{5cm}
%		\begin{picture}(1,1)
%		\put(0,0){\line(0,1){1}}
%		\put(0,0){\line(1,0){1}}
%		\put(0,0){\line(1,1){1}}
%		\put(0,0){\line(1,2){.5}}
%		\put(0,0){\line(1,3){.3333}}
%		\put(0,0){\line(1,4){.25}}
%		\put(0,0){\line(1,5){.2}}
%		\put(0,0){\line(1,6){.1667}}
%		\put(0,0){\line(2,1){1}}
%		\put(0,0){\line(2,3){.6667}}
%		\put(0,0){\line(2,5){.4}}
%		\put(0,0){\line(3,1){1}}
%		\put(0,0){\line(3,2){1}}
%		\put(0,0){\line(3,4){.75}}
%		\put(0,0){\line(3,5){.6}}
%		\put(0,0){\line(4,1){1}}
%		\put(0,0){\line(4,3){1}}
%		\put(0,0){\line(4,5){.8}}
%		\put(0,0){\line(5,1){1}}
%		\put(0,0){\line(5,2){1}}
%		\put(0,0){\line(5,3){1}}
%		\put(0,0){\line(5,4){1}}
%		\put(0,0){\line(5,6){.8333}}
%		\put(0,0){\line(6,1){1}}
%		\put(0,0){\line(6,5){1}}
%		\end{picture}
%	\end{center}
%	\caption{\label{fig_circulo}A delimitação do espaço}
%	\legend{Fonte: os autores}
%\end{figure}

%Ou então figuras podem ser incorporadas de arquivos externos, como é o caso da \autoref{fig_grafico}. Se a figura que ser incluída se tratar de um diagrama, um gráfico ou uma ilustração que você mesmo produza, priorize o uso de imagens vetoriais no formato PDF. Com isso, o tamanho do arquivo final do trabalho será menor, e as imagens terão uma apresentação melhor, principalmente quando impressas, uma vez que imagens vetorias são perfeitamente escaláveis para qualquer dimensão. Nesse caso, se for utilizar o Microsoft Excel para produzir gráficos, ou o Microsoft Word para produzir ilustrações, exporte-os como PDF e os incorpore ao documento conforme o exemplo abaixo. No entanto, para manter a coerência no uso de software livre (já que você está usando \LaTeX e \abnTeX), teste a ferramenta \textsf{InkScape}\index{InkScape} (\url{http://inkscape.org/}). Ela é uma excelente opção de código-livre para produzir ilustrações vetoriais, similar ao CorelDraw\index{CorelDraw} ou ao Adobe Illustrator\index{Adobe Illustrator}. De todo modo, caso não seja possível utilizar arquivos de imagens como PDF, utilize qualquer outro formato, como JPEG, GIF, BMP, etc. Nesse caso, você pode tentar aprimorar as imagens incorporadas com o software livre \textsf{Gimp}\index{Gimp} (\url{http://www.gimp.org/}). Ele é uma alternativa livre ao Adobe Photoshop\index{Adobe Photoshop}.

%\begin{figure}[htb]
% 	\begin{center}
% 	    \includegraphics[scale=0.5]{01.jpg}
% 	\end{center}
% 	\caption{\label{fig_grafico}Gráfico produzido em Excel e salvo como PDF}
% 	\legend{Fonte: \citeonline[p. 24]{araujo2012}}
% \end{figure}

% ----------------------------------------------
%\section{Enumerações: alíneas e subalíneas}
% ----------------------------------------------

%\index{alíneas}\index{subalíneas}\index{incisos}Quando for necessário enumerar
%os diversos assuntos de uma seção que não possua título, esta deve ser
%subdividida em alíneas \cite[4.2]{NBR6024:2012}:

%\begin{alineas}

%  \item os diversos assuntos que não possuam título próprio, dentro de uma mesma seção, devem ser subdivididos em alíneas\footnote{As notas devem ser digitadas ou datilografadas dentro das margens, ficando separadas do texto por um espaço simples de entre as linhas e por filete de 5 cm, a partir da margem esquerda. Devem ser alinhadas, a partir da segunda linha da mesma nota, abaixo da primeira letra da primeira palavra, de forma a destacar o expoente, sem espaço entre elas e com fonte menor. \citeonline[5.2.1]{NBR14724:2011}}; 
%  \item o texto que antecede as alíneas termina em dois pontos;
%  \item as alíneas devem ser indicadas alfabeticamente, em letra minúscula, seguida de parêntese. Utilizam-se letras dobradas, quando esgotadas as letras do alfabeto;
%  \item as letras indicativas das alíneas devem apresentar recuo em relação à margem esquerda;
%  \item o texto da alínea deve começar por letra minúscula e terminar em ponto-e-vírgula, exceto a última alínea que termina em ponto final;
%  \item o texto da alínea deve terminar em dois pontos, se houver subalínea;
 % \item a segunda e as seguintes linhas do texto da alínea começa sob a primeira letra do texto da própria alínea;
 % \item subalíneas \cite[4.3]{NBR6024:2012} devem ser conforme as alíneas a seguir:

% \begin{alineas}
%     \item as subalíneas devem começar por travessão seguido de espaço;
%     \item as subalíneas devem apresentar recuo em relação à alínea;
%     \item o texto da subalínea deve começar por letra minúscula e terminar em ponto-e-vírgula. A última subalínea deve terminar em ponto final, se não houver alínea subsequente;
%     \item a segunda e as seguintes linhas do texto da subalínea começam sob a primeira letra do texto da própria subalínea.
%  \end{alineas}
  
%  \item no \abnTeX\ estão disponíveis os ambientes \texttt{incisos} e
%  \texttt{subalineas}, que em suma são o mesmo que se criar outro nível de \texttt{alineas}, como nos exemplos à seguir:
  
%  \begin{incisos}
%    \item \textit{Um novo inciso em itálico};
%  \end{incisos}
  
%  \item Alínea em \textbf{negrito}:
  
%  \begin{subalineas}
%    \item \textit{Uma subalínea em itálico};
%    \item \underline{\textit{Uma subalínea em itálico e sublinhado}}; 
%  \end{subalineas}
  
%  \item Última alínea com \emph{ênfase}.
  
%\end{alineas}


% ----------------------------------------------
%\section{Inclução de outros arquivos}\label{sec-include}
% ----------------------------------------------
%É uma boa prática dividir o seu documento em diversos arquivos, e não apenas escrever tudo em um único. Esse recurso foi utilizado neste documento. Para incluir diferentes arquivos em um arquivo principal, de modo que cada arquivo incluído fique em uma página diferente, utilize o comando:
%\begin{verbatim}
%   \include{documento-a-ser-incluido}      % sem a extensão .tex
%\end{verbatim}

%Para incluir documentos sem quebra de páginas, utilize:
%\begin{verbatim}
 %  \input{documento-a-ser-incluido}      % sem a extensão .tex
%\end{verbatim}

% ||||||||||||||||||||||||||||||||||||||||||||||
% CAP 3 - MATERIAIS, FERRAMENTAS E MÉTODOS
% ||||||||||||||||||||||||||||||||||||||||||||||
\chapter{Materiais, ferramentas e métodos}

\section{Proposta de pesquisa}

\lipsum[2-3]

\begin{figure}[htp]
	\centering
	\caption{\label{fig:inrush-fig02} Logo Unisinos.}
	\includegraphics[width = 0.8\linewidth]{images/unisinos.png}
	\legend{Fonte: Adaptado de \citeonline{NBR14724:2011}.}
\end{figure}

% ||||||||||||||||||||||||||||||||||||||||||||||
% CAP 4 - CRONOGRAMA
% ||||||||||||||||||||||||||||||||||||||||||||||
\chapter{Cronograma} \label{cronograma}

A Tabela \ref{tab_cronograma} apresenta o cronograma de execução das atividades desta proposta.\\

\begin{table}[!htpb]
\centering
\caption{Cronograma das atividades previstas}
\setlength{\tabcolsep}{5pt} 
\begin{tabular}{|c|c|c|c|c|c|c|c|c|c|c|c|c|c|}\hline
 & \multicolumn{12}{c|}{Meses}\\ \cline{2-13}
\raisebox{1.5ex}{Etapa} & 1 & 2 & 3 & 4 & 5 & 6 & 7 & 8 & 9 & 10 & 11 & 12 \\ \hline \hline
A1 & \cellcolor{black} &   &   &   &   &   &   &   &   &   &   &   \\ \hline
A2 &   & \cellcolor{black} &   &   &   &   &   &   &   &   &   &   \\ \hline
A3 &   & \cellcolor{black} & \cellcolor{black} &   &   &   &   &   &   &   &   &   \\ \hline
A4 &   &   &   & \cellcolor{black} & \cellcolor{black} &   &   &   &   &   &   &   \\ \hline
A5 &   &   &   &   &   & \cellcolor{black} & \cellcolor{black} &   &   &   &   &   \\ \hline
A6 &   &   &   &   &   &   &   & \cellcolor{black} & \cellcolor{black} &   &   &   \\ \hline
A7 &   &   &   &   &   &   &   &   &   & \cellcolor{black} & \cellcolor{black} & \cellcolor{black} \\ \hline
\end{tabular} 
\legend{Fonte: Produzido pelo autor  }
\label{tab_cronograma}
\end{table} 

% ||||||||||||||||||||||||||||||||||||||||||||||
% CONCLUSÃO (capítulo sem numeração e presente no sumário)
% ||||||||||||||||||||||||||||||||||||||||||||||
%\chapter*[Conclusão]{Conclusão}
%\addcontentsline{toc}{chapter}{Conclusão}
% Utilize caso a conclusão seja um capitulo sem numeracao.

% ||||||||||||||||||||||||||||||||||||||||||||||
% CAPITULO 5
% ||||||||||||||||||||||||||||||||||||||||||||||
% \chapter{Conclusão}

% \lipsum[11-12]


% \section{Trabalhos futuros}

% \lipsum[21-22]

% ----------------------------------------------
% Finaliza a parte no bookmark do PDF para que se inicie o bookmark na raiz e adiciona espaço de parte no Sumário
% ----------------------------------------------
\phantompart

% Arquivo de elementos pós-textuais: referências, apêndices e anexos 
\input{3-postextual}

% ----------------------------------------------
\end{document}
% ----------------------------------------------

% ||||||||||||||||||||||||||||||||||||||||||||||
% ALGUNS EXEMPLOS DE CODIGO LATEX
% ||||||||||||||||||||||||||||||||||||||||||||||

% ----------------------------------------------
%	Figure example
% ----------------------------------------------
% \begin{figure}[htp]
% \centering
% \caption{\label{fig:x} Aqui vai o caption da imagem.}
% \includegraphics[width = 0.9\linewidth ]{images/x.png}
% \legend{Fonte: Adaptado de \citeonline{x}.}
% \end{figure}

% ----------------------------------------------
%	Figure example
% ----------------------------------------------
%\begin{figure}[htp]
%\centering
%\includegraphics[width = 0.5\linewidth ]{dlayer.jpg}
%\caption{\label{fig:1}Complanar waveguide parameters using two dielectric layers \cite{cwccs}.} 
%\end{figure}

% ----------------------------------------------
%	Equation example \ref{eqn:1}
% ----------------------------------------------
%\begin{eqnarray}
%\epsilon _{r} & = & \epsilon _{r}'(1-i\tan(\delta))
%\label{eqn:1}
%\end{eqnarray}

% ----------------------------------------------
%	Complex equations example \ref{eqn:2}
% ----------------------------------------------
%\begin{eqnarray}
%C_{1} & = & 2\epsilon _{0}(\epsilon _{r1}-1)\frac{K(k_{1})}{K(k_{1}')}\nonumber\\
%& = & 2\epsilon _{0}(\epsilon _{r1}'-i\tan(\delta)\epsilon _{r1}'-1)\frac{K(k_{1})}{K(k_{1}')}\nonumber\\
%& = & 2\epsilon _{0}(\epsilon _{r1}'-1)\frac{K(k_{1})}{K(k_{1}')}-i2\epsilon _{0}(\tan(\delta)\epsilon _{r1}')\frac{K(k_{1})}{K(k_{1}')}\\
%C_{2} & = & 2\epsilon _{0}(\epsilon _{r2}-\epsilon _{r1})\frac{K(k_{2})}{K(k_{2}')}\nonumber\\
%& = & 2\epsilon _{0}(\epsilon _{r2}'-i\tan(\delta)\epsilon _{r2}'-\epsilon _{r1}'+i\tan(\delta)\epsilon _{r1}')\frac{K(k_{2})}{K(k_{2}')}\nonumber\\
%& = & 2\epsilon _{0}(\epsilon _{r2}'-\epsilon _{r1}')\frac{K(k_{2})}{K(k_{2}')}-i2\epsilon _{0}(\tan(\delta)\epsilon _{r2}'-\tan(\delta)\epsilon _{r1}')\frac{K(k_{2})}{K(k_{2}')}\\
%C_{vac} & = & 4\epsilon _{0}\frac{K(k_{0})}{K(k_{0}')}
%\label{eqn:2}
%\end{eqnarray}

% ----------------------------------------------
%	Table example \ref{tab:1} 
% ----------------------------------------------
%\begin{table}[htb]
%\caption{Constants and Parameters.}
%\begin{center}
%\begin{tabular}{|c|c|c|c|}
%\hline
%\bfseries CONSTANT & \bfseries VALUE & \bfseries CONSTANT & \bfseries VALUE \\
%\hline \hline
%$\epsilon _{0}$ & 8.8540$\times$10$^{-12}$ & $c$ & 299792458~m/s \\
%\hline
%$\epsilon _{r1}$ & 11.7 & $h_{1}$ & 300~$\mu$m \\
%\hline
%$\epsilon _{r2}$ & 7.5 & $h_{2}$ & 200~nm \\
%\hline
%$\delta _{1}$ & 10$^{-4}$ & $\delta _{2}$ & 10$^{-2}$ $\sim$ 10$^{-3}$ \\
%\hline
%\end{tabular}
%\end{center}
%\label{tab:1}
%\end{table}

% ----------------------------------------------
%	Complex table example \ref{tab:2}
% ----------------------------------------------
%\begin{table}[htb]
%\caption{Results and dimensions.}
%\begin{center}
%\begin{tabular}{|c|c|c|c|c|}
%\hline
%\bfseries $f_{0}$ [MHz] & \bfseries $S$ [$\mu$m] & \bfseries $W$ [$\mu$m] & \bfseries $d$ [cm] & \bfseries $C_{c}$ [fF] \\
%\hline
%\hline
%\multirow{4}{*}{650} & \multirow { 2}{*}{2} & \multirow{ 2}{*}{1} & \multirow{ 2}{*}{9.4966} & 10 \\
%\cline{5-5}
%& & & & 30 \\
%\cline{2-4} \cline{5-5}
%& \multirow { 2}{*}{4.5} & \multirow { 2}{*}{2.3} & \multirow { 2}{*}{9.3155} &         10 \\
%\cline{5-5}
%& & & & 30 \\
%\hline
%\multirow { 4}{*}{6000} &  \multirow { 2}{*}{2} &  \multirow { 2}{*}{1} &  \multirow { 2}{*}{1.0288} &          1 \\
%\cline{5-5}
%& & & & 3 \\
%\cline{2-4} \cline{5-5}
%& \multirow { 2}{*}{4.5} &  \multirow { 2}{*}{2.3} &  \multirow { 2}{*}{1.0092} &          1 \\
%\cline{5-5}
%& & & & 3 \\
%\hline

%\end{tabular}
%\end{center}
%\label{tab:2}
%\end{table}

% ----------------------------------------------
%	Two graphics in one \ref{fig:2}
% ----------------------------------------------
%\begin{figure}[htp]
%\centering
%\begin{tabular}{cc}
%(a) & (b) \\
%\includegraphics[width = 0.5\linewidth ]{DM_6G_4p5_3_6.jpg} &
%\includegraphics[width = 0.5\linewidth ]{DM_650_4p5_10_6.jpg}
%\end{tabular}
%\caption{captions} 
%\label{fig:2}
%\end{figure}

% ----------------------------------------------
%	Four graphics in one table \ref{fig:3}
% ----------------------------------------------
%\begin{figure}[htp]
%\centering
%\begin{tabular}{cc}
%(a) & (b) \\
%\includegraphics[width = 0.5\linewidth ]{DM_650_4p5_10_1.jpg} &
%\includegraphics[width = 0.5\linewidth ]{DM_650_4p5_30_1.jpg} \\
%(c) & (d) \\
%\includegraphics[width = 0.5\linewidth ]{DM_6G_4p5_1_1A.jpg} &
%\includegraphics[width = 0.5\linewidth ]{DM_6G_4p5_3_1.jpg} \\
%\end{tabular}
%\caption{Captions} 
%\label{fig:3}
%\end{figure}

% ----------------------------------------------
%	Quote and footnote
% ----------------------------------------------
%\begin{quote} ``adkfjahsldkfjashflkasdjfhadslkfjhasdlfkjadshflsda
%kjdshflkasjdfhalskfjhadslfksdajhfladskjfhsda
%kajsdfhlaksdjfhasdkl'' \footnote{test footnote}
%\end{quote}

% ----------------------------------------------
%	PDF Annotation
% ----------------------------------------------
%\pdfannot % generic annotation
%width 10cm % the dimension of the annotation can be controlled
%height 0cm % via <rule spec>; if some of dimensions in
%depth 4cm % <rule spec> is not given, the corresponding
% value of the parent box will be used.
%{ %
%/Subtype /Text % text annotation
%/Open true % if given then the text annotation will be opened
%/Contents % text contents
%(write comments in here...)
%}%
%

% ----------------------------------------------
% Tutorial para compilação offline
% ----------------------------------------------
% Faça o download dos seguintes softwares: Miktex (basic) e texStudio
% https://miktex.org/download
% http://www.texstudio.org/
% Instale o Miktex e selecione para instalar os pacotes automaticamente... Install Packages on the Fly = Yes
% Depois de instalado abra o "Miktex Package Manager (Admin)" e já instale os pacotes "abntex2" e o "cm-super".
% Instale o texstudio
% Faça o download dos arquivos do projeto no Overleaf "Download as ZIP"
% Descompacte no diretório de projeto desejado... use um diretório para cada projeto pois durante a compilação serão gerados diversos arquivos
% Abra o "main.tex" do projeto no texstudio e clique no icone >> para compilar o projeto... tecla de atalho = F5.
